% (c) Nikita Lisitsa, lisyarus@gmail.com, 2025

\documentclass[10pt]{beamer}

\usepackage[T2A]{fontenc}
\usepackage[russian]{babel}
\usepackage{minted}

\usepackage[most]{tcolorbox}
\tcbuselibrary{minted}

\usepackage{graphicx}
\graphicspath{ {./images/} }

\usetheme{metropolis}

\usepackage{adjustbox}

\usepackage{color}
\usepackage{soul}

\usepackage{hyperref}

\definecolor{codebg}{RGB}{29,35,49}

\setminted{fontsize=\footnotesize}

\makeatletter
\newcommand{\slideimage}[1]{
  \begin{figure}
    \begin{adjustbox}{width=\textwidth, totalheight=\textheight-2\baselineskip-2\baselineskip,keepaspectratio}
      \includegraphics{#1}
    \end{adjustbox}
  \end{figure}
}
\makeatother

\title{Компьютерная графика}
\subtitle{Практика 1: Рисуем треугольник}
\date{2026}

\setbeamertemplate{footline}[frame number]

\usemintedstyle{solarized-light}

\begin{document}

\frame{\titlepage}

\begin{frame}[fragile]
\frametitle{}
\textbf{О практиках}
\begin{itemize}
\item Задания идут \textit{последовательно}, полагаясь на предыдущие задания этой же практики
\item Иногда задания \textit{перекрывают} предыдущие задания, тогда нужно \textit{закомментировать код предыдущего задания}, чтобы его тоже можно было проверить
\item Некоторые задания занимают \textit{больше одного слайда} :)
\end{itemize}
\end{frame}

\begin{frame}[fragile]
\frametitle{}
\textbf{О практиках}
\begin{itemize}
\item Вся инициализация спрятана от вас в мини-библиотеку \nolinkurl{wgpu_app}
\item В C++ из неё вам будут нужны \mintinline{cpp}|app.device()|, \mintinline{cpp}|app.queue()|, \mintinline{cpp}|app.surface_format()|
\item В Rust из неё вам будут нужны \mintinline{cpp}|gpu.device| и \mintinline{cpp}|gpu.queue|, \mintinline{cpp}|gpu.surface_format|
\item Для C++ в репозитории есть маленькая библиотека математики -- её использовать можно \textbf{только с третьей практики} :)
\end{itemize}
\end{frame}

\begin{frame}[fragile]
\frametitle{}
\textbf{О практиках}
\begin{itemize}
\item Мы постоянно будем загружать файлы с шейдерами, 3D моделями, текстурами, сценами
\item Они будут лежать в корне конкретной практики
\item Для удобства в C++ уже есть константа \mintinline{cpp}|projectRoot|, в Rust -- \mintinline{rust}|PROJECT_ROOT|
\item В C++ есть вспомогательная функция \mintinline{cpp}|loadFile| в \mintinline{cpp}|<file_utils.hpp>|
\end{itemize}
\end{frame}

\begin{frame}[fragile]
\frametitle{}
\textbf{Работа с WebGPU}
\begin{itemize}
\item Функции создания WebGPU-объектов всегда принимают т.н. \textit{descriptor} -- небольшую структуру с описанием объекта
\item В C++ вместо них можно передавать \mintinline{cpp}|nullptr|, если вам подойдёт дефолтная настройка
\item На C++ их удобно инициализировать специальными макросами, например
\usemintedstyle{lightbulb}
\begin{tcblisting}{listing engine=minted, minted language=cpp, minted options={fontsize=\scriptsize}, listing only, colback=codebg, colframe=codebg, rounded corners}
WGPURenderPassColorAttachment colorAttachment =
    WGPU_RENDER_PASS_COLOR_ATTACHMENT_INIT;
\end{tcblisting}
\item На Rust можно использовать
\begin{tcblisting}{listing engine=minted, minted language=rust, minted options={fontsize=\scriptsize}, listing only, colback=codebg, colframe=codebg, rounded corners}
wgpu::RenderPassColorAttachment {
    // описание полей...
    ..Default::default()
}
\end{tcblisting}
\end{itemize}
\end{frame}

\begin{frame}[fragile]
\frametitle{}
\textbf{Работа с WebGPU}
\begin{itemize}
\item Так как \nolinkurl{wgpu-native} -- API для языка C, массивы в дескрипторах задаются как количество элементов + указатель на первый элемент, например
\usemintedstyle{lightbulb}
\begin{tcblisting}{listing engine=minted, minted language=cpp, minted options={fontsize=\scriptsize}, listing only, colback=codebg, colframe=codebg, rounded corners}
struct WGPUFragmentState {
    // ...
    size_t targetCount;
    WGPUColorTargetState const * targets;
};
\end{tcblisting}
\end{itemize}
\end{frame}

\begin{frame}[fragile]
\frametitle{}
\textbf{Работа с WebGPU}
\begin{itemize}
\item В C++ все созданные объекты нужно удалить, когда они не нужны, иначе будут утечки памяти
\pause
\item Это делается функциями вида \mintinline{cpp}|wgpuCommandBufferRelease|
\pause
\item На самом деле \nolinkurl{wgpu} держит счётчик ссылок на объекты, так что делать \textit{release} не страшно, даже если \nolinkurl{wgpu} внутри всё ещё использует объект
\pause
\item В Rust \nolinkurl{wgpu} делает это автоматически (через Drop trait)
\end{itemize}
\end{frame}

\begin{frame}[fragile]
\frametitle{Задание 1}
\begin{itemize}
\item Создаём \textbf{command encoder} -- объект, позволяющий записать последовательность команд для выполнения на GPU
\item Превращаем его в \textbf{command buffer} -- собственно последовательность команд
\item Отправляем их в \textbf{очередь команд} для исполнения -- пока ни одной команды нет, так что ничего интересного не произойдёт
\item Всё это в основном цикле рендеринга, сразу после \mintinline{rust}|surface_texture.create_view|.
\begin{itemize}
\item C++: \mintinline{cpp}|wgpuDeviceCreateCommandEncoder()|, \mintinline{cpp}|wgpuCommandEncoderFinish()|, \mintinline{cpp}|wgpuQueueSubmit()|, не забываем про release
\item Rust: \mintinline{rust}|device.create_command_encoder()|, \mintinline{rust}|encoder.finish()|, \mintinline{rust}|queue.submit()|
\end{itemize}
\end{itemize}
\end{frame}

\begin{frame}[fragile]
\frametitle{Задание 2}
\begin{itemize}
\item Через command encoder создаём \textbf{render pass}, который очистит наш экран фиксированным цветом
\item Привязываем к нему \textbf{color attachment}, привязанный к нашему экрану (C++: \mintinline{cpp}|target|, Rust: \mintinline{rust}|view|)
\item \mintinline{cpp}|StoreOp| -- \mintinline{cpp}|StoreOp|, \mintinline{cpp}|LoadOp| -- \mintinline{cpp}|Clear|, цвет выберите сами
\item Всё это \textbf{до} \mintinline{rust}|encoder.finish()|
\begin{itemize}
\item C++: \mintinline{cpp}|WGPURenderPassColorAttachment|, \mintinline{cpp}|WGPURenderPassDescriptor|, \mintinline{cpp}|wgpuCommandEncoderBeginRenderPass()|, \mintinline{cpp}|wgpuRenderPassEncoderEnd()|
\item Rust: \mintinline{rust}|encoder.begin_render_pass()|
\end{itemize}
\end{itemize}
\end{frame}

\begin{frame}[fragile]
\frametitle{Задание 2}
\slideimage{task_2.png}
\end{frame}

\begin{frame}[fragile]
\frametitle{Задание 3}
\begin{itemize}
\item Пишем наш первый шейдер: создаём файл \nolinkurl{shader.wgsl} с магическим содержанием
\usemintedstyle{lightbulb}
\begin{tcblisting}{listing engine=minted, minted language=rust, minted options={fontsize=\scriptsize}, listing only, colback=codebg, colframe=codebg, rounded corners}
@vertex
fn vertexMain(@builtin(vertex_index) vertexIndex: u32)
    -> @builtin(position) vec4f 
{
    var positions = array<vec2f, 3>(
        vec2f( 0.0,  0.5),
        vec2f(-0.5, -0.5),
        vec2f( 0.5, -0.5),
    );
    return vec4f(positions[vertexIndex], 0.0, 1.0);
}

@fragment
fn fragmentMain() -> @location(0) vec4f {
    return vec4f(0.95, 0.35, 0.15, 1.0);
}
\end{tcblisting}
\usemintedstyle{solarized-light}
\item (цвет во \mintinline{rust}|fragmentMain| можно выбрать любой)
\end{itemize}
\end{frame}

\begin{frame}[fragile]
\frametitle{Задание 4}
\begin{itemize}
\item При инициализации программы (до \mintinline{cpp}|while (running)| в C++, \mintinline{cpp}|WgpuApp::new| в Rust) загружаем наш файл с кодом шейдера и создаём \textbf{shader module}
\item Он просто скомпилируется, но пока не используется -- попробуйте сделать в нём ошибку (например, добавить лишнюю запятую) и увидеть, что компиляция падает с ошибкой
\item C++: \mintinline{cpp}|WGPUShaderSourceWGSL|, \mintinline{cpp}|WGPUShaderModuleDescriptor|, \mintinline{cpp}|wgpuDeviceCreateShaderModule()|
\item Rust: \mintinline{rust}|device.create_shader_module()|, \mintinline{rust}|wgpu::ShaderSource::Wgsl|
\end{itemize}
\end{frame}

\begin{frame}[fragile]
\frametitle{Задание 5}
\begin{itemize}
\item Создаём \textbf{render pipeline} -- основной объект, описывающий, что, как, и в каком формате рисовать
\item Важно заполнить:
\usemintedstyle{lightbulb}
\begin{tcblisting}{listing engine=minted, minted language=cpp, minted options={fontsize=\scriptsize}, listing only, colback=codebg, colframe=codebg, rounded corners}
VertexState
    module = созданный нами shader module
    entry_point = "vertexMain"
FragmentState
    module = созданный нами shader module
    entry_point = "fragmentMain"
    targets = [{
        format = app.surfaceFormat()
        write_mask = ALL
    }]
PrimitiveState
    topology = TriangleList
\end{tcblisting}
\usemintedstyle{solarized-light}
\item Остальное можно оставить дефолтами
\end{itemize}
\end{frame}

\begin{frame}[fragile]
\frametitle{Задание 5}
\begin{itemize}
\item Пока мы не используем этот pipeline, так что ничего не произойдёт
\item C++: \mintinline{cpp}|WGPUColorTargetState|, \mintinline{cpp}|WGPUFragmentState|, \mintinline{cpp}|WGPURenderPipelineDescriptor|, \mintinline{cpp}|wgpuDeviceCreateRenderPipeline()|
\item Rust: \mintinline{rust}|device.create_render_pipeline()|
\end{itemize}
\end{frame}

\begin{frame}[fragile]
\frametitle{Задание 6}
\begin{itemize}
\item Используем созданный \textbf{render pipeline}: перед \mintinline{rust}|encoder.end()| устанавливаем наш pipeline и вызываем рисование одного треугольника
\item \nolinkurl{vertex_count = 3}, \nolinkurl{instance_count = 1}, \nolinkurl{first_vertex} и \nolinkurl{first_instance} по нулям
\item Должен появиться цветной треугольник :)
\item C++: \mintinline{cpp}|wgpuRenderPassEncoderSetPipeline|, \mintinline{cpp}|wgpuRenderPassEncoderDraw()|
\item Rust: \mintinline{rust}|encoder.set_pipeline()|, \mintinline{rust}|encoder.draw()|
\end{itemize}
\end{frame}

\begin{frame}[fragile]
\frametitle{Задание 6}
\slideimage{task_6.png}
\end{frame}

\begin{frame}[fragile]
\frametitle{Задание 7*}
\begin{itemize}
\item Раскрашиваем треугольник в шахматном порядке
\item Нужно как-то изменить \mintinline{rust}|fragmentMain| в коде шейдера
\item Полезные переменные и функции WGSL: \mintinline{rust}|@builtin(position)|, \mintinline{rust}|modf()|
\end{itemize}
\end{frame}

\begin{frame}[fragile]
\frametitle{Задание 7*}
\slideimage{task_7.png}
\end{frame}

\end{document}